\chapter{Introduction to Statistical Learning}
\label{ch:topic1}

\section{Regression Function}
The function
\[
m(x):= \text{E}[Y \mid X = x]
\]
is called the \textit{regression function}. It describes how the \textbf{expected value of Y changes with X}.

\begin{table}[h]
\begin{tabular}{@{}l p{3cm} l@{}}
\toprule
Possible cuts & Type & Interpretation\\
\midrule
$\text{E}[Y \mid X = x]$ & number & conditional mean at $x$\\
$m(\cdot)$ & function & maps $x$ to the conditional mean\\
$\text{E}[Y\mid X] = m(X)$ & random variable & $m$ evaluated at random $X$\\
\botrule
\end{tabular}
\end{table}

\section{Key Properties}
\subsection{Conditional Expectation}
Conditional expectation behaves much like ordinary expectation, but with $X$ treated as fixed/unknown.

\begin{itemize}
    \item \textbf{Linearity: } $\text{E}[aY + bZ \mid X]=a\text{E}[Y \mid X]+b\text{E}[Z \mid X]$
    \item \textbf{Taking out what is known: } $\text{E}[g(X)Y \mid X] = g(X)\text{E}[Y \mid X]$
    \item \textbf{Independence:} If $X$ and $Y$ are independent, then $\text{E}[Y \mid X] = \text{E}[Y]$
    \item \textbf{Tower property:} $\text{E}(\text{E}[Y \mid X])=\text{E}[Y]$ and $\text{E}(\text{E}[Y \mid Z, X] \mid X) =\text{E}(Y \mid X)$
\end{itemize}

\subsection{Conditional Variance}
The conditional mean
\[
m(X)=\text{E}[Y \mid X]
\]
describes the part of $Y$ that can be predicted from $X$.
The variation that remains after observing $X$ is measured by
\[
\operatorname{Var}(X \mid Y)=\text{E}[\{{Y-\text{E}[Y \mid X]\}^2 \mid X}]
\]

\begin{definition}[Law of Total Variance]
\begin{equation}
\operatorname{Var}(Y)=\operatorname{E}[\operatorname{Var}(Y \mid X)] + \operatorname{Var}(\operatorname{E}[Y \mid X])
\end{equation}
\end{definition}
where $\operatorname{Var}(Y)$ is the total variation, $\operatorname{E}[\operatorname{Var}(Y \mid X)]$ is the variation remaining after observing $X$, and $\operatorname{Var}(\operatorname{E}[Y \mid X])$ is the variation explained by $X$.

\begin{example}[Toy Example]
Let $X \in \{0,1,2\}$ be the number of exclamation marks in an email, and let
\begin{equation*}
Y=\left\{\begin{array}{@{}c@{}c}
1,&\quad \mbox{spam}\\
0,&\quad \mbox{genuine}
\end{array}
\right.
\end{equation*}
Suppose the joint distribution of $(X,Y)$ is
\begin{table}[h]
    \centering
    \begin{tabular}{ccccc}\toprule
         &$x=0$& $x=1$&$x=2$& $p_Y(y)$\\\midrule
         $y=0$&   0.4& 0.15&0.05& 0.60\\
         $y=1$&   0.05& 0.10&0.25& 0.40\\
 $p_X(x)$& 0.45& 0.25& 0.30&1\\ \bottomrule 
    \end{tabular}
    \label{tab:placeholder}
\end{table}

From the marginals,
\begin{equation}
    \operatorname{E}[X] =0.85, \quad \operatorname{E}[Y]=P(\text{spam)}=0.40 
\end{equation}
\end{example}





    